% !TeX TS-program = xelatex
\documentclass[10pt, a4paper]{article}

% ── Geometry ──
\usepackage[
  top=0.4in,
  bottom=0.4in,
  left=0.6in,
  right=0.6in
]{geometry}

% ── Fonts (Latin) ──
\usepackage{fontspec}

% ── Packages ──
\usepackage{enumitem}
\usepackage{titlesec}
\usepackage{hyperref}
\usepackage{xcolor}

% ── Colors ──
\definecolor{linkblue}{HTML}{0563C1}
\definecolor{sectioncolor}{HTML}{000000}
\definecolor{rulecolor}{HTML}{333333}

% ── Hyperlinks ──
\hypersetup{
  colorlinks=true,
  urlcolor=linkblue,
  linkcolor=linkblue,
  pdfauthor={Tianshuo Han},
  pdftitle={Tianshuo Han - Resume},
}

% ── No page numbers ──
\pagestyle{empty}

% ── Section style: bold with underline ──
\titleformat{\section}
  {\normalsize\bfseries\color{sectioncolor}}
  {}{0em}{}
  [\vspace{-0.6em}\textcolor{rulecolor}{\rule{\textwidth}{0.8pt}}]
\titlespacing{\section}{0pt}{8pt}{4pt}

% ── List style ──
\setlist[itemize]{
  leftmargin=1.5em,
  itemsep=0pt,
  parsep=0pt,
  topsep=1pt,
}

% ── Helper commands ──

\newcommand{\resumeentry}[2]{%
  \noindent\textbf{#1}\hfill{#2}\par
}

\newcommand{\resumesub}[1]{%
  \noindent{#1}\par\vspace{1pt}
}

\newcommand{\pub}[3]{%
  \noindent\textbf{#1}\par
  \noindent{#2}\par
  \noindent{#3}\par\vspace{1pt}
}

% ── Reduce spacing ──
\setlength{\parindent}{0pt}
\setlength{\parskip}{0pt}

% ══════════════════════════════════════════
\begin{document}

% ── Header ──
\begin{center}
  {\Large\bfseries Tianshuo Han}\\[4pt]
  (+86)\,173-3953-9195 \enspace$\mid$\enspace
  \href{mailto:hantianshuo@iie.ac.cn}{hantianshuo@iie.ac.cn} \enspace$\mid$\enspace
  \href{https://github.com/keymaker-arch}{github.com/keymaker-arch}
\end{center}

\vspace{-2pt}

% ── Education ──
\section{Education}

\resumeentry{University of Chinese Academy of Sciences \enspace|\enspace Ph.D. Candidate}{Sep 2021 --- Jan 2027 (Expected)}
\resumesub{Institute of Information Engineering \enspace|\enspace Direct Ph.D. Program}

Focused on \textbf{Linux kernel} and \textbf{operating systems} research, with solid theoretical foundation and extensive hands-on experience in kernel internals, binary analysis, and system-level software testing. Research published at top-tier conferences including IEEE S\&P, USENIX, and ACM CCS.

\vspace{6pt}

\resumeentry{China Agricultural University \enspace|\enspace Undergraduate}{Sep 2017 --- Jun 2021}
\resumesub{Biological Sciences, College of Life Sciences \enspace|\enspace Bachelor of Science}

Outstanding academic performance with multiple academic scholarships. Awarded National Second Prize in the China Undergraduate Mathematical Contest in Modeling and Gold Medal in the iGEM Competition.

% ── Publications ──
\section{Publications}

\pub{Kernel Fuzzing of Distributed File Systems via Intra-Cluster Protocol Poisoning}
{[First Author], \textbf{Tianshuo Han}, et al. (2nd author)}
{ACM CCS, 2026 (Under Submission), CCF-A}

\begin{itemize}
  \item Proposed intra-cluster protocol poisoning as a novel kernel attack surface for distributed file systems. Built a fuzzer that interposes on inter-node traffic and mutates protocol messages in flight, guided by static-analysis-based contamination-scope quantification and protocol-state-synchronized progressive exploration.
  \item Discovered XX previously unknown kernel bugs (heap overflows, use-after-free, null-pointer dereferences) across five distributed file systems running the latest Linux kernel.
\end{itemize}

\vspace{4pt}

\pub{NetPanic: A Fuzzer-in-the-Middle Approach to Auditing the Linux Kernel Network Stack}
{\textbf{Tianshuo Han}, Zong Cao, Zhen Dong, Xiapu Luo, Zhenyu Song, and Jian Liu}
{47th IEEE S\&P, 2026, CCF-A}

\begin{itemize}
  \item Conducted the first systematic fuzzing-based analysis of the Linux kernel network stack. Proposed the Fuzzer-in-the-Middle architecture, achieving over 100x throughput and 400\% code coverage improvement over the state-of-the-art, discovering 15 previously unknown kernel bugs.
  \item Independently completed the overall architecture design and implementation. Led the paper writing effort.
\end{itemize}

\vspace{4pt}

\pub{CARDSHARK: Understanding and Stabilizing Linux Kernel Concurrency Bugs Against the Odds}
{\textbf{Tianshuo Han}, Xiaorui Gong, and Jian Liu}
{33rd USENIX Symposium, 2024, CCF-A}

\begin{itemize}
  \item Conducted the first in-depth analysis of the root causes of instability in triggering kernel concurrency bugs, and proposed a general-purpose technique that significantly improves the reproducibility and stability of concurrency bug triggering.
  \item Independently completed the mechanism analysis, technique development, and performance evaluation. Led the paper writing effort.
\end{itemize}

\vspace{4pt}

\pub{Reviving Discarded Bugs: Analyzing Previously Unanalyzable Linux Kernel Bugs Through Control Metadata Fields}
{Hao Zhang, Jie Lu, Shaomin Chen, \textbf{Tianshuo Han}, Bolun Zhang, Jian Liu, Xiaorui Gong}
{32nd ACM CCS, 2025, CCF-A}

\begin{itemize}
  \item Proposed a novel technique based on manipulating critical metadata fields in kernel structures, enhancing analysis capabilities in constrained scenarios and expanding the scope of analyzable kernel bugs.
  \item Deeply involved in defining the core thesis and responding to reviewer comments. Participated in evaluation experiment design and CodeQL development.
\end{itemize}

\vspace{4pt}

\pub{Whole-genome sequence and comparative analysis of Trichoderma asperellum ND-1 reveal its unique enzymatic system for efficient biomass degradation}
{Fengzhen Zheng, \textbf{Tianshuo Han}, Abdul Basit, Junquan Liu, Ting Miao, and Wei Jiang}
{Catalysts 12(4), 2022, IF=4.16}

\begin{itemize}
  \item Revealed the unique lignocellulose degradation enzyme system of \textit{Trichoderma asperellum} ND-1 through whole-genome sequencing and secretome analysis. Responsible for genomic data analysis and visualization.
\end{itemize}

\end{document}
